\chapter{主要行业的职业卫生}
\setcounter{chapter}{12}
\chapterintro{
\textbf{【教学大纲要求】}

\imp{了解：}采矿、冶金、化工、建筑、农业等行业的主要职业有害因素及防护要点。
}

\section{采矿行业}

\begin{itemize}[leftmargin=*]
    \item \imp{主要危害因素}：粉尘（矽尘、煤尘等）→尘肺；噪声和振动（凿岩机、爆破）；高温高湿（深矿井）；有毒气体（爆破产生CO、NOx等）；不良体位
    \item \imp{防护要点}：湿式作业、通风除尘、噪声控制、个体防护
\end{itemize}

\section{冶金行业}

\begin{itemize}[leftmargin=*]
    \item \imp{主要危害因素}：高温和热辐射→中暑；粉尘（金属粉尘、煤尘）；有毒气体（CO、SO$_2$、氟化物等）；噪声
    \item \imp{防护要点}：隔热降温、通风排毒、噪声控制
\end{itemize}

\section{化工行业}

\begin{itemize}[leftmargin=*]
    \item \imp{主要危害因素}：各类化学毒物（有机溶剂、刺激性气体、致癌物等）；火灾爆炸危险
    \item \imp{防护要点}：密闭化自动化生产、通风排毒、个人防护、应急救援预案
\end{itemize}

\section{建筑行业}

\begin{itemize}[leftmargin=*]
    \item \imp{主要危害因素}：高处坠落、物体打击、粉尘（水泥尘、矽尘）、噪声和振动、高温和寒冷
    \item \imp{防护要点}：安全防护设施、个体防护、安全教育培训
\end{itemize}

\section{农业}

\begin{itemize}[leftmargin=*]
    \item \imp{主要危害因素}：农药中毒、有机粉尘（农民肺）、人畜共患传染病（布氏杆菌病、炭疽）、不良气象条件、机械伤害
    \item \imp{防护要点}：农药安全使用、通风防尘、个人防护、健康教育
\end{itemize}

\newpage

% =====================================================================
%  附录：核心名词解释速记表
% =====================================================================
\chapter*{附录：核心名词解释速记表}
\addcontentsline{toc}{chapter}{附录：核心名词解释速记表}

\begin{table}[htbp]
\centering
\begin{tabularx}{\textwidth}{lX}
\hline
\textbf{名词} & \textbf{核心释义} \\
\hline
职业卫生 & 研究劳动条件对健康的影响，识别、评价、预测、控制不良劳动条件 \\
职业医学 & 对受职业危害因素损害的个体进行检测、诊断、治疗和康复 \\
职业病 & 职业活动中接触有害因素引起、具有医学和法律双重含义的疾病 \\
工作有关疾病 & 职业因素为多因素之一的疾病，改善工作条件可控制或缓解 \\
三级预防 & 一级（病因预防）、二级（发病预防——健康监护）、三级（康复预防）\\
职业中毒 & 工作人员在生产过程中接触生产性毒物引起的中毒 \\
中毒性肺水肿 & 吸入高浓度刺激性气体后肺泡及间质过量体液潴留 \\
铅中毒 & 卟啉代谢障碍→贫血、周围神经病（垂腕）、铅绞痛、铅线 \\
汞中毒 & 易兴奋症（特有）、意向性震颤、口腔炎 \\
锰中毒 & 帕金森氏综合征（面具脸、慌张步态、齿轮样肌张力、震颤）\\
苯中毒 & CNS麻醉（急性）、造血系统损害→再障→白血病（慢性，1类致癌物）\\
有机磷中毒 & 不可逆抑制AChE→ACh蓄积→M样+N样+CNS症状。解毒：阿托品+氯磷定 \\
CO中毒 & 与Hb结合→HbCO（樱桃红色），亲和力比O$_2$大200--300倍 \\
HCN中毒 & 抑制细胞色素氧化酶Fe$^{3+}$→中断呼吸链→细胞窒息。解毒：亚硝酸钠-硫代硫酸钠 \\
尘肺 & 长期吸入矿物性粉尘所致的以肺弥漫性纤维化为主的疾病 \\
矽肺 & 吸入游离SiO$_2$粉尘所致，矽结节为特征性病理改变 \\
石棉肺 & 吸入石棉纤维所致，弥漫性间质纤维化、石棉小体、胸膜斑 \\
矽肺结核 & 矽肺合并肺结核，最常见最严重并发症 \\
石棉致肺癌/间皮瘤 & 石棉IARC 1类致癌物，吸烟有协同致癌作用 \\
防尘八字方针 & 革、水、密、风、护、管、教、查 \\
\hline
\end{tabularx}
\end{table}

\begin{table}[htbp]
\centering
\begin{tabularx}{\textwidth}{lX}
\hline
\textbf{名词} & \textbf{核心释义} \\
\hline
热射病 & 最严重中暑类型：高热>40$^\circ$C、无汗、意识障碍 \\
减压病 & 高气压下减压过快氮气释放形成气泡栓塞 \\
噪声聋 & 长期接触生产性噪声引起的以高频听力下降为主的感音性耳聋 \\
健康工人效应（HWE） & 工人总死亡率低于一般人群的现象，职业流行病学重要偏倚来源 \\
热适应 & 机体在长期高温环境下产生的形态、结构和生理功能的适应性变化 \\
热习服 & 短期反复接触高温后生理功能补偿性调整，7--14天形成，脱离后消退 \\
射频电磁场 & 高频（HF）、超高频（VHF/UHF）、微波（MW），热效应和非热效应 \\
电光性眼炎 & 紫外线引起的急性角膜结膜炎，潜伏期4--12h \\
振动性白指（VWF） & 寒冷诱发手指苍白（雷诺现象），局部振动病最特征性改变 \\
职业性肿瘤 & 职业活动中接触致癌因素引起的肿瘤，潜伏期长 \\
确定致癌物 & 流行病学证据充分+IARC 1类，流行病学调查为最重要识别手段 \\
职业接触限值 & MAC（最高容许浓度）、PC-TWA（时间加权）、PC-STEL（短时间）\\
生物监测 & 测定生物材料中化学物或其代谢产物以评价内暴露水平 \\
职业健康监护 & 就业前体检+定期体检+离岗时体检+应急检查 \\
\hline
\end{tabularx}
\end{table}

\newpage

\vspace*{3cm}

\end{document}
\newpage